% !TeX program = xelatex
\documentclass[12pt,a4paper]{article}
\usepackage[margin=2cm]{geometry}
\usepackage{fontspec}
\usepackage{amsmath,amssymb,enumitem}
\setmainfont{Times New Roman}
\setlength{\parindent}{0pt}
\setlength{\parskip}{6pt}
\begin{document}
\begin{center}
{\Large\bfseries Đề thi Toán vào lớp 10 Hà Nội năm 2020 -- phần Đại số}
\end{center}

\textbf{Bài I.} Cho hai biểu thức
\[
A=\frac{\sqrt{x}+1}{\sqrt{x}+2},\qquad
B=\frac{3}{\sqrt{x}-1}-\frac{\sqrt{x}+5}{x-1},
\quad x\ge 0,\ x\ne 1.
\]
\begin{enumerate}[label=\arabic*)]
  \item Tính giá trị của biểu thức $A$ khi $x=4$.
  \item Chứng minh $B=\dfrac{2}{\sqrt{x}+1}$.
  \item Tìm tất cả giá trị của $x$ để biểu thức $P=2A\cdot B+\sqrt{x}$ đạt giá trị nhỏ nhất.
\end{enumerate}

\textbf{Bài II.}
\begin{enumerate}[label=\arabic*)]
  \item Giải bài toán sau bằng cách lập phương trình hoặc hệ phương trình:

  Quãng đường từ nhà An đến nhà Bình dài $3$ km. Buổi sáng, An đi bộ từ nhà An đến nhà Bình. Buổi chiều cùng ngày, An đi xe đạp từ nhà Bình về nhà An trên cùng quãng đường đó với vận tốc lớn hơn vận tốc đi bộ của An là $9$ km/h. Tính vận tốc đi bộ của An, biết thời gian đi buổi chiều ít hơn thời gian đi buổi sáng là $45$ phút. (Giả định rằng An đi bộ với vận tốc không đổi trên toàn bộ quãng đường đó.)
  \item Một quả bóng bàn có dạng một hình cầu có bán kính bằng $2$ cm. Tính diện tích bề mặt của quả bóng bàn đó (lấy $\pi\approx 3{,}14$).
\end{enumerate}

\textbf{Bài III.}
\begin{enumerate}[label=\arabic*)]
  \item Giải hệ phương trình
  \[
  \begin{cases}
  2x+\dfrac{3}{y-1}=5,\\[4pt]
  4x-\dfrac{1}{y-1}=3.
  \end{cases}
  \]
  \item Trong mặt phẳng tọa độ $Oxy$, xét đường thẳng $(d):y=mx+4$ với $m\ne 0$.
  \begin{enumerate}[label=\alph*)]
    \item Gọi $A$ là giao điểm của đường thẳng $(d)$ và trục $Oy$. Tìm tọa độ điểm $A$.
    \item Tìm tất cả giá trị của $m$ để đường thẳng $(d)$ cắt trục $Ox$ tại điểm $B$ sao cho tam giác $OAB$ là tam giác cân.
  \end{enumerate}
\end{enumerate}

\textbf{Bài V.} Giải phương trình
\[
\sqrt{x}+\sqrt{3x-2}=x^2+1.
\]
\end{document}
